\documentclass[10pt,a4paper]{amsart}
\usepackage{amsmath,amssymb,amsthm,graphicx}
% For commutative diagrams we use Xy-pic
\usepackage[all]{xy}

% We use 2cell for 2-commutative diagrams.
\xyoption{2cell}
% Theorem environments.
%
\theoremstyle{plain}
\newtheorem{theorem}[subsection]{Theorem}
\newtheorem{proposition}[subsection]{Proposition}
\newtheorem{lemma}[subsection]{Lemma}

\theoremstyle{definition}
\newtheorem{definition}[subsection]{Definition}
\newtheorem{example}[subsection]{Example}
\newtheorem{exercise}[subsection]{Exercise}
\newtheorem{situation}[subsection]{Situation}

\theoremstyle{remark}
\newtheorem{remark}[subsection]{Remark}
\newtheorem{remarks}[subsection]{Remarks}

\numberwithin{equation}{subsection}

% Macros
%
\def\lim{\mathop{\mathrm{lim}}\nolimits}
\def\colim{\mathop{\mathrm{colim}}\nolimits}
\def\Spec{\mathop{\mathrm{Spec}}}
\def\Hom{\mathop{\mathrm{Hom}}\nolimits}
\def\Ext{\mathop{\mathrm{Ext}}\nolimits}
\def\SheafHom{\mathop{\mathcal{H}\!\mathit{om}}\nolimits}
\def\SheafExt{\mathop{\mathcal{E}\!\mathit{xt}}\nolimits}
\def\Sch{\mathit{Sch}}
\def\Mor{\mathop{\mathrm{Mor}}\nolimits}
\def\Ob{\mathop{\mathrm{Ob}}\nolimits}
\def\Sh{\mathop{\mathit{Sh}}\nolimits}
\def\NL{\mathop{N\!L}\nolimits}
\def\CH{\mathop{\mathrm{CH}}\nolimits}
\def\proetale{{pro\text{-}\acute{e}tale}}
\def\etale{{\acute{e}tale}}
\def\QCoh{\mathit{QCoh}}
\def\Ker{\mathop{\mathrm{Ker}}}
\def\Im{\mathop{\mathrm{Im}}}
\def\Coker{\mathop{\mathrm{Coker}}}
\def\Coim{\mathop{\mathrm{Coim}}}

% Boxtimes
%
\DeclareMathSymbol{\boxtimes}{\mathbin}{AMSa}{"02}

%
% Macros for moduli stacks/spaces
%
\def\QCohstack{\mathcal{QC}\!\mathit{oh}}
\def\Cohstack{\mathcal{C}\!\mathit{oh}}
\def\Spacesstack{\mathcal{S}\!\mathit{paces}}
\def\Quotfunctor{\mathrm{Quot}}
\def\Hilbfunctor{\mathrm{Hilb}}
\def\Curvesstack{\mathcal{C}\!\mathit{urves}}
\def\Polarizedstack{\mathcal{P}\!\mathit{olarized}}
\def\Complexesstack{\mathcal{C}\!\mathit{omplexes}}
% \Pic is the operator that assigns to X its picard group, usage \Pic(X)
% \Picardstack_{X/B} denotes the Picard stack of X over B
% \Picardfunctor_{X/B} denotes the Picard functor of X over B
\def\Pic{\mathop{\mathrm{Pic}}\nolimits}
\def\Picardstack{\mathcal{P}\!\mathit{ic}}
\def\Picardfunctor{\mathrm{Pic}}
\def\Deformationcategory{\mathcal{D}\!\mathit{ef}}
\newenvironment{reference}{}{}
\newenvironment{slogan}{}{}
\newenvironment{history}{}{}
\usepackage{cleveref}
\crefname{originallemma}{original lemma}{original lemmas}
\Crefname{originallemma}{Original lemma}{Original lemmas}
\makeatletter
\newlabel{lemma-U-finite-above-x}{{03JS}{source}{}{sourcepointer-1}{}}
\newlabel{lemma-U-finite-above-x@cref}{{[originallemma][0][]03JS}{[1][0][]source}{}{}{}}
\makeatother
\crefname{originallemma}{original lemma}{original lemmas}
\Crefname{originallemma}{Original lemma}{Original lemmas}
\makeatletter
\newlabel{lemma-R-finite-above-x}{{03JU}{source}{}{sourcepointer-2}{}}
\newlabel{lemma-R-finite-above-x@cref}{{[originallemma][0][]03JU}{[1][0][]source}{}{}{}}
\makeatother
\crefname{originallemma}{original lemma}{original lemmas}
\Crefname{originallemma}{Original lemma}{Original lemmas}
\makeatletter
\newlabel{lemma-UR-finite-above-x}{{03JV}{source}{}{sourcepointer-3}{}}
\newlabel{lemma-UR-finite-above-x@cref}{{[originallemma][0][]03JV}{[1][0][]source}{}{}{}}
\makeatother
\crefname{originallemma}{original lemma}{original lemmas}
\Crefname{originallemma}{Original lemma}{Original lemmas}
\makeatletter
\newlabel{lemma-U-universally-bounded}{{03JT}{source}{}{sourcepointer-4}{}}
\newlabel{lemma-U-universally-bounded@cref}{{[originallemma][0][]03JT}{[1][0][]source}{}{}{}}
\makeatother
\crefname{originallemma}{original lemma}{original lemmas}
\Crefname{originallemma}{Original lemma}{Original lemmas}
\makeatletter
\newlabel{lemma-characterize-very-reasonable}{{03IH}{source}{}{sourcepointer-5}{}}
\newlabel{lemma-characterize-very-reasonable@cref}{{[originallemma][0][]03IH}{[1][0][]source}{}{}{}}
\makeatother
\crefname{originallemma}{original lemma}{original lemmas}
\Crefname{originallemma}{Original lemma}{Original lemmas}
\makeatletter
\newlabel{spaces-properties-lemma-quasi-separated-quasi-compact-pieces}{{03W7}{source}{}{sourcepointer-6}{}}
\newlabel{spaces-properties-lemma-quasi-separated-quasi-compact-pieces@cref}{{[originallemma][0][]03W7}{[1][0][]source}{}{}{}}
\makeatother
\crefname{originallemma}{original lemma}{original lemmas}
\Crefname{originallemma}{Original lemma}{Original lemmas}
\makeatletter
\newlabel{morphisms-lemma-locally-quasi-finite-qc-source-universally-bounded}{{03JA}{source}{}{sourcepointer-7}{}}
\newlabel{morphisms-lemma-locally-quasi-finite-qc-source-universally-bounded@cref}{{[originallemma][0][]03JA}{[1][0][]source}{}{}{}}
\makeatother
\crefname{originallemma}{original lemma}{original lemmas}
\Crefname{originallemma}{Original lemma}{Original lemmas}
\makeatletter
\newlabel{lemma-descent-universally-bounded}{{03JO}{source}{}{sourcepointer-8}{}}
\newlabel{lemma-descent-universally-bounded@cref}{{[originallemma][0][]03JO}{[1][0][]source}{}{}{}}
\makeatother
\crefname{originallemma}{original lemma}{original lemmas}
\Crefname{originallemma}{Original lemma}{Original lemmas}
\makeatletter
\newlabel{lemma-base-change-universally-bounded}{{03JN}{source}{}{sourcepointer-9}{}}
\newlabel{lemma-base-change-universally-bounded@cref}{{[originallemma][0][]03JN}{[1][0][]source}{}{}{}}
\makeatother
\crefname{originallemma}{original lemma}{original lemmas}
\Crefname{originallemma}{Original lemma}{Original lemmas}
\makeatletter
\newlabel{lemma-composition-universally-bounded}{{03JM}{source}{}{sourcepointer-10}{}}
\newlabel{lemma-composition-universally-bounded@cref}{{[originallemma][0][]03JM}{[1][0][]source}{}{}{}}
\makeatother
\crefname{originallemma}{original lemma}{original lemmas}
\Crefname{originallemma}{Original lemma}{Original lemmas}
\makeatletter
\newlabel{lemma-universally-bounded-finite-fibres}{{03JQ}{source}{}{sourcepointer-11}{}}
\newlabel{lemma-universally-bounded-finite-fibres@cref}{{[originallemma][0][]03JQ}{[1][0][]source}{}{}{}}
\makeatother
\crefname{originallemma}{original lemma}{original lemmas}
\Crefname{originallemma}{Original lemma}{Original lemmas}
\makeatletter
\newlabel{lemma-there-is-a-scheme-integral-over}{{09YB}{source}{}{sourcepointer-12}{}}
\newlabel{lemma-there-is-a-scheme-integral-over@cref}{{[originallemma][0][]09YB}{[1][0][]source}{}{}{}}
\makeatother
\crefname{originalproposition}{original proposition}{original propositions}
\Crefname{originalproposition}{Original proposition}{Original propositions}
\makeatletter
\newlabel{spaces-properties-proposition-locally-quasi-separated-open-dense-scheme}{{06NH}{source}{}{sourcepointer-13}{}}
\newlabel{spaces-properties-proposition-locally-quasi-separated-open-dense-scheme@cref}{{[originalproposition][0][]06NH}{[1][0][]source}{}{}{}}
\makeatother
\crefname{originallemma}{original lemma}{original lemmas}
\Crefname{originallemma}{Original lemma}{Original lemmas}
\makeatletter
\newlabel{morphisms-lemma-generically-finite}{{02NW}{source}{}{sourcepointer-14}{}}
\newlabel{morphisms-lemma-generically-finite@cref}{{[originallemma][0][]02NW}{[1][0][]source}{}{}{}}
\makeatother
\crefname{originallemma}{original lemma}{original lemmas}
\Crefname{originallemma}{Original lemma}{Original lemmas}
\makeatletter
\newlabel{spaces-morphisms-lemma-quasi-finite-separated-quasi-affine}{{0ABS}{source}{}{sourcepointer-15}{}}
\newlabel{spaces-morphisms-lemma-quasi-finite-separated-quasi-affine@cref}{{[originallemma][0][]0ABS}{[1][0][]source}{}{}{}}
\makeatother
\crefname{originallemma}{original lemma}{original lemmas}
\Crefname{originallemma}{Original lemma}{Original lemmas}
\makeatletter
\newlabel{spaces-morphisms-lemma-properties-normalization}{{0ABP}{source}{}{sourcepointer-16}{}}
\newlabel{spaces-morphisms-lemma-properties-normalization@cref}{{[originallemma][0][]0ABP}{[1][0][]source}{}{}{}}
\makeatother
\crefname{originallemma}{original lemma}{original lemmas}
\Crefname{originallemma}{Original lemma}{Original lemmas}
\makeatletter
\newlabel{spaces-morphisms-lemma-normalization-in-integral}{{0825}{source}{}{sourcepointer-17}{}}
\newlabel{spaces-morphisms-lemma-normalization-in-integral@cref}{{[originallemma][0][]0825}{[1][0][]source}{}{}{}}
\makeatother
\crefname{originallemma}{original lemma}{original lemmas}
\Crefname{originallemma}{Original lemma}{Original lemmas}
\makeatletter
\newlabel{spaces-morphisms-lemma-characterize-normalization}{{0823}{source}{}{sourcepointer-18}{}}
\newlabel{spaces-morphisms-lemma-characterize-normalization@cref}{{[originallemma][0][]0823}{[1][0][]source}{}{}{}}
\makeatother
\crefname{originalsection}{original section}{original sections}
\Crefname{originalsection}{Original section}{Original sections}
\makeatletter
\newlabel{spaces-morphisms-section-scheme-theoretic-closure}{{0831}{source}{}{sourcepointer-19}{}}
\newlabel{spaces-morphisms-section-scheme-theoretic-closure@cref}{{[originalsection][0][]0831}{[1][0][]source}{}{}{}}
\makeatother
\crefname{originallemma}{original lemma}{original lemmas}
\Crefname{originallemma}{Original lemma}{Original lemmas}
\makeatletter
\newlabel{spaces-morphisms-lemma-quasi-compact-immersion}{{088G}{source}{}{sourcepointer-20}{}}
\newlabel{spaces-morphisms-lemma-quasi-compact-immersion@cref}{{[originallemma][0][]088G}{[1][0][]source}{}{}{}}
\makeatother
\crefname{originallemma}{original lemma}{original lemmas}
\Crefname{originallemma}{Original lemma}{Original lemmas}
\makeatletter
\newlabel{spaces-morphisms-lemma-semi-diagonal}{{03KO}{source}{}{sourcepointer-21}{}}
\newlabel{spaces-morphisms-lemma-semi-diagonal@cref}{{[originallemma][0][]03KO}{[1][0][]source}{}{}{}}
\makeatother
\crefname{originallemma}{original lemma}{original lemmas}
\Crefname{originallemma}{Original lemma}{Original lemmas}
\makeatletter
\newlabel{lemma-quasi-compact-reasonable-stratification}{{07S8}{source}{}{sourcepointer-22}{}}
\newlabel{lemma-quasi-compact-reasonable-stratification@cref}{{[originallemma][0][]07S8}{[1][0][]source}{}{}{}}
\makeatother
\crefname{originallemma}{original lemma}{original lemmas}
\Crefname{originallemma}{Original lemma}{Original lemmas}
\makeatletter
\newlabel{spaces-morphisms-lemma-composition-integral}{{03ZR}{source}{}{sourcepointer-23}{}}
\newlabel{spaces-morphisms-lemma-composition-integral@cref}{{[originallemma][0][]03ZR}{[1][0][]source}{}{}{}}
\makeatother
\title{Integral scheme covers generically finite-etale and locally refining a prescribed etale cover}
\author{The Stacks Project Authors}
\date{stacks-project:decent-spaces.tex:local-snapshot-20261011; GNU FDL 1.2 or later}
\begin{document}
\maketitle
\noindent\textit{Editorial note.} At thed=1 base step the author has already shownX is a scheme. The integral identity coverY=X satisfies the conclusion and locally factors through the open covering; the printedY=V there is an immediate variable-choice slip because a general open cover need not be integral. Original text retained.\par\medskip
\noindent\textit{Editorial note.} The phrase finite overV\_j in the normalization base-change paragraph can be read as integral/affine, which is all that the scheme and quasi-compactness conclusion uses. No finiteness of normalization over an arbitrary base is added.\par\medskip
\noindent\textit{Allowed standard prerequisites (editorial).}\begin{itemize}
\item Explicit prerequisite boundary. Quasi-compact quasi-separated algebraic spaces, dense open scheme loci, relative normalization, quasi-finite separated maps, etale diagonal sections and fibre-cardinality stratification. The bounded-fibre argument is retained author support.
\end{itemize}\medskip
\ifdefined\hypertarget\hypertarget{sourcepointer-1}{}\fi
\noindent\textit{Source reference.} Original lemma 03JS: focused-seeds/AGP/sources/stacks-snapshot-20261011/tags.txt, source lines 4650--4650. Stable Stacks tag for an expressly allowed foundational or notational input; the counted construction remains the two concrete core obligations. The retained local core is not replaced by this pointer. See standard\_prerequisites.\par\medskip
\ifdefined\hypertarget\hypertarget{sourcepointer-2}{}\fi
\noindent\textit{Source reference.} Original lemma 03JU: focused-seeds/AGP/sources/stacks-snapshot-20261011/tags.txt, source lines 4652--4652. Stable Stacks tag for an expressly allowed foundational or notational input; the counted construction remains the two concrete core obligations. The retained local core is not replaced by this pointer. See standard\_prerequisites.\par\medskip
\ifdefined\hypertarget\hypertarget{sourcepointer-3}{}\fi
\noindent\textit{Source reference.} Original lemma 03JV: focused-seeds/AGP/sources/stacks-snapshot-20261011/tags.txt, source lines 4653--4653. Stable Stacks tag for an expressly allowed foundational or notational input; the counted construction remains the two concrete core obligations. The retained local core is not replaced by this pointer. See standard\_prerequisites.\par\medskip
\ifdefined\hypertarget\hypertarget{sourcepointer-4}{}\fi
\noindent\textit{Source reference.} Original lemma 03JT: focused-seeds/AGP/sources/stacks-snapshot-20261011/tags.txt, source lines 4651--4651. Stable Stacks tag for an expressly allowed foundational or notational input; the counted construction remains the two concrete core obligations. The retained local core is not replaced by this pointer. See standard\_prerequisites.\par\medskip
\ifdefined\hypertarget\hypertarget{sourcepointer-5}{}\fi
\noindent\textit{Source reference.} Original lemma 03IH: focused-seeds/AGP/sources/stacks-snapshot-20261011/tags.txt, source lines 4603--4603. Stable Stacks tag for an expressly allowed foundational or notational input; the counted construction remains the two concrete core obligations. The retained local core is not replaced by this pointer. See standard\_prerequisites.\par\medskip
\ifdefined\hypertarget\hypertarget{sourcepointer-6}{}\fi
\noindent\textit{Source reference.} Original lemma 03W7: focused-seeds/AGP/sources/stacks-snapshot-20261011/tags.txt, source lines 5102--5102. Stable Stacks tag for an expressly allowed foundational or notational input; the counted construction remains the two concrete core obligations. The retained local core is not replaced by this pointer. See standard\_prerequisites.\par\medskip
\ifdefined\hypertarget\hypertarget{sourcepointer-7}{}\fi
\noindent\textit{Source reference.} Original lemma 03JA: focused-seeds/AGP/sources/stacks-snapshot-20261011/tags.txt, source lines 4632--4632. Stable Stacks tag for an expressly allowed foundational or notational input; the counted construction remains the two concrete core obligations. The retained local core is not replaced by this pointer. See standard\_prerequisites.\par\medskip
\ifdefined\hypertarget\hypertarget{sourcepointer-8}{}\fi
\noindent\textit{Source reference.} Original lemma 03JO: focused-seeds/AGP/sources/stacks-snapshot-20261011/tags.txt, source lines 4646--4646. Stable Stacks tag for an expressly allowed foundational or notational input; the counted construction remains the two concrete core obligations. The retained local core is not replaced by this pointer. See standard\_prerequisites.\par\medskip
\ifdefined\hypertarget\hypertarget{sourcepointer-9}{}\fi
\noindent\textit{Source reference.} Original lemma 03JN: focused-seeds/AGP/sources/stacks-snapshot-20261011/tags.txt, source lines 4645--4645. Stable Stacks tag for an expressly allowed foundational or notational input; the counted construction remains the two concrete core obligations. The retained local core is not replaced by this pointer. See standard\_prerequisites.\par\medskip
\ifdefined\hypertarget\hypertarget{sourcepointer-10}{}\fi
\noindent\textit{Source reference.} Original lemma 03JM: focused-seeds/AGP/sources/stacks-snapshot-20261011/tags.txt, source lines 4644--4644. Stable Stacks tag for an expressly allowed foundational or notational input; the counted construction remains the two concrete core obligations. The retained local core is not replaced by this pointer. See standard\_prerequisites.\par\medskip
\ifdefined\hypertarget\hypertarget{sourcepointer-11}{}\fi
\noindent\textit{Source reference.} Original lemma 03JQ: focused-seeds/AGP/sources/stacks-snapshot-20261011/tags.txt, source lines 4648--4648. Stable Stacks tag for an expressly allowed foundational or notational input; the counted construction remains the two concrete core obligations. The retained local core is not replaced by this pointer. See standard\_prerequisites.\par\medskip
\ifdefined\hypertarget\hypertarget{sourcepointer-12}{}\fi
\noindent\textit{Source reference.} Original lemma 09YB: focused-seeds/AGP/sources/stacks-snapshot-20261011/tags.txt, source lines 12559--12559. Stable Stacks tag for an expressly allowed foundational or notational input; the counted construction remains the two concrete core obligations. The retained local core is not replaced by this pointer. See standard\_prerequisites.\par\medskip
\ifdefined\hypertarget\hypertarget{sourcepointer-13}{}\fi
\noindent\textit{Source reference.} Original proposition 06NH: focused-seeds/AGP/sources/stacks-snapshot-20261011/tags.txt, source lines 8536--8536. Stable Stacks tag for an expressly allowed foundational or notational input; the counted construction remains the two concrete core obligations. The retained local core is not replaced by this pointer. See standard\_prerequisites.\par\medskip
\ifdefined\hypertarget\hypertarget{sourcepointer-14}{}\fi
\noindent\textit{Source reference.} Original lemma 02NW: focused-seeds/AGP/sources/stacks-snapshot-20261011/tags.txt, source lines 3495--3495. Stable Stacks tag for an expressly allowed foundational or notational input; the counted construction remains the two concrete core obligations. The retained local core is not replaced by this pointer. See standard\_prerequisites.\par\medskip
\ifdefined\hypertarget\hypertarget{sourcepointer-15}{}\fi
\noindent\textit{Source reference.} Original lemma 0ABS: focused-seeds/AGP/sources/stacks-snapshot-20261011/tags.txt, source lines 13030--13030. Stable Stacks tag for an expressly allowed foundational or notational input; the counted construction remains the two concrete core obligations. The retained local core is not replaced by this pointer. See standard\_prerequisites.\par\medskip
\ifdefined\hypertarget\hypertarget{sourcepointer-16}{}\fi
\noindent\textit{Source reference.} Original lemma 0ABP: focused-seeds/AGP/sources/stacks-snapshot-20261011/tags.txt, source lines 13027--13027. Stable Stacks tag for an expressly allowed foundational or notational input; the counted construction remains the two concrete core obligations. The retained local core is not replaced by this pointer. See standard\_prerequisites.\par\medskip
\ifdefined\hypertarget\hypertarget{sourcepointer-17}{}\fi
\noindent\textit{Source reference.} Original lemma 0825: focused-seeds/AGP/sources/stacks-snapshot-20261011/tags.txt, source lines 10240--10240. Stable Stacks tag for an expressly allowed foundational or notational input; the counted construction remains the two concrete core obligations. The retained local core is not replaced by this pointer. See standard\_prerequisites.\par\medskip
\ifdefined\hypertarget\hypertarget{sourcepointer-18}{}\fi
\noindent\textit{Source reference.} Original lemma 0823: focused-seeds/AGP/sources/stacks-snapshot-20261011/tags.txt, source lines 10238--10238. Stable Stacks tag for an expressly allowed foundational or notational input; the counted construction remains the two concrete core obligations. The retained local core is not replaced by this pointer. See standard\_prerequisites.\par\medskip
\ifdefined\hypertarget\hypertarget{sourcepointer-19}{}\fi
\noindent\textit{Source reference.} Original section 0831: focused-seeds/AGP/sources/stacks-snapshot-20261011/tags.txt, source lines 10271--10271. Stable Stacks tag for an expressly allowed foundational or notational input; the counted construction remains the two concrete core obligations. The retained local core is not replaced by this pointer. See standard\_prerequisites.\par\medskip
\ifdefined\hypertarget\hypertarget{sourcepointer-20}{}\fi
\noindent\textit{Source reference.} Original lemma 088G: focused-seeds/AGP/sources/stacks-snapshot-20261011/tags.txt, source lines 10461--10461. Stable Stacks tag for an expressly allowed foundational or notational input; the counted construction remains the two concrete core obligations. The retained local core is not replaced by this pointer. See standard\_prerequisites.\par\medskip
\ifdefined\hypertarget\hypertarget{sourcepointer-21}{}\fi
\noindent\textit{Source reference.} Original lemma 03KO: focused-seeds/AGP/sources/stacks-snapshot-20261011/tags.txt, source lines 4684--4684. Stable Stacks tag for an expressly allowed foundational or notational input; the counted construction remains the two concrete core obligations. The retained local core is not replaced by this pointer. See standard\_prerequisites.\par\medskip
\ifdefined\hypertarget\hypertarget{sourcepointer-22}{}\fi
\noindent\textit{Source reference.} Original lemma 07S8: focused-seeds/AGP/sources/stacks-snapshot-20261011/tags.txt, source lines 9893--9893. Stable Stacks tag for an expressly allowed foundational or notational input; the counted construction remains the two concrete core obligations. The retained local core is not replaced by this pointer. See standard\_prerequisites.\par\medskip
\ifdefined\hypertarget\hypertarget{sourcepointer-23}{}\fi
\noindent\textit{Source reference.} Original lemma 03ZR: focused-seeds/AGP/sources/stacks-snapshot-20261011/tags.txt, source lines 5230--5230. Stable Stacks tag for an expressly allowed foundational or notational input; the counted construction remains the two concrete core obligations. The retained local core is not replaced by this pointer. See standard\_prerequisites.\par\medskip
\begin{lemma}
\label{lemma-there-is-a-scheme-integral-over-refined}
Let $S$ be a scheme. Let $X$ be a quasi-compact and quasi-separated
algebraic space over $S$ such that $|X|$ has finitely many irreducible
components.
\begin{enumerate}
\item There exists a surjective integral morphism $Y \to X$ where $Y$
is a scheme such that $f$ is finite \'etale over a quasi-compact
dense open $U \subset X$,
\item given a surjective \'etale morphism $V \to X$ we may choose
$Y \to X$ such that for every $y \in Y$ there is an open neighbourhood
$W \subset Y$ such that $W \to X$ factors through $V$.
\end{enumerate}
\end{lemma}

\begin{lemma}
\label{lemma-bounded-fibres}
Let $S$ be a scheme. Let $X$ be an algebraic space over $S$.
Consider the following conditions on $X$:
\begin{itemize}
\item[] $(\alpha)$ For every $x \in |X|$, the equivalent conditions of
Lemma \ref{lemma-U-finite-above-x}
hold.
\item[] $(\beta)$ For every $x \in |X|$, the equivalent conditions of
Lemma \ref{lemma-R-finite-above-x}
hold.
\item[] $(\gamma)$ For every $x \in |X|$, the equivalent conditions of
Lemma \ref{lemma-UR-finite-above-x}
hold.
\item[] $(\delta)$ The equivalent conditions of
Lemma \ref{lemma-U-universally-bounded}
hold.
\item[] $(\epsilon)$ The equivalent conditions of
Lemma \ref{lemma-characterize-very-reasonable}
hold.
\item[] $(\zeta)$ The space $X$ is Zariski locally quasi-separated.
\item[] $(\eta)$ The space $X$ is quasi-separated
\item[] $(\theta)$ The space $X$ is representable, i.e., $X$ is a scheme.
\item[] $(\iota)$ The space $X$ is a quasi-separated scheme.
\end{itemize}
We have
$$
\xymatrix{
& (\theta) \ar@{=>}[rd] & & & &  \\
(\iota) \ar@{=>}[ru] \ar@{=>}[rd] & &
(\zeta) \ar@{=>}[r] &
(\epsilon) \ar@{=>}[r] &
(\delta) \ar@{=>}[r] &
(\gamma) \ar@{<=>}[r] & (\alpha) + (\beta) \\
& (\eta) \ar@{=>}[ru] & & & &
}
$$
\end{lemma}

\begin{proof}
The implication $(\gamma) \Leftrightarrow (\alpha) + (\beta)$ is immediate.
The implications in the diamond on the left are clear from the
definitions.

\medskip\noindent
Assume $(\zeta)$, i.e., that $X$ is Zariski locally quasi-separated.
Then $(\epsilon)$ holds by
Properties of Spaces, Lemma
\ref{spaces-properties-lemma-quasi-separated-quasi-compact-pieces}.

\medskip\noindent
Assume $(\epsilon)$. By
Lemma \ref{lemma-characterize-very-reasonable}
there exists
a Zariski open covering $X = \bigcup X_i$ such that for each $i$
there exists a scheme $U_i$ and a quasi-compact surjective \'etale morphism
$U_i \to X_i$. Choose an $i$ and an affine open subscheme $W \subset U_i$.
It suffices to show that $W \to X$ has universally bounded fibres, since then
the family of all these morphisms $W \to X$ covers $X$.
To do this we consider the diagram
$$
\xymatrix{
W \times_X U_i \ar[r]_-p \ar[d]_q & U_i \ar[d] \\
W \ar[r] & X
}
$$
Since $W \to X$ factors through $X_i$ we see that
$W \times_X U_i = W \times_{X_i} U_i$, and hence $q$ is quasi-compact.
Since $W$ is affine this implies that the scheme $W \times_X U_i$
is quasi-compact. Thus we may apply
Morphisms, Lemma
\ref{morphisms-lemma-locally-quasi-finite-qc-source-universally-bounded}
and we conclude that $p$ has universally bounded fibres. From
Lemma \ref{lemma-descent-universally-bounded}
we conclude that $W \to X$ has universally bounded fibres as well.

\medskip\noindent
Assume $(\delta)$. Let $U$ be an affine scheme, and let $U \to X$ be an \'etale
morphism. By assumption the fibres of the morphism $U \to X$ are universally
bounded. Thus also the fibres of both projections $R = U \times_X U \to U$
are universally bounded, see
Lemma \ref{lemma-base-change-universally-bounded}.
And by
Lemma \ref{lemma-composition-universally-bounded}
also the fibres of $R \to X$ are universally bounded.
Hence for any $x \in X$ the fibres of $|U| \to |X|$ and $|R| \to |X|$
over $x$ are finite, see
Lemma \ref{lemma-universally-bounded-finite-fibres}.
In other words, the equivalent conditions of
Lemma \ref{lemma-UR-finite-above-x}
hold. This proves that $(\delta) \Rightarrow (\gamma)$.
\end{proof}

\begin{proof}
The proof is the (roughly) same as the proof of
Lemma \ref{lemma-there-is-a-scheme-integral-over}
with additional technical comments to obtain
the dense quasi-compact open $U$ (and unfortunately
changes in notation to keep track of $U$).

\medskip\noindent
Part (1) is the special case of part (2) where $V = X$.

\medskip\noindent
Proof of (2). Choose a surjective \'etale morphism $V' \to V$
where $V'$ is a scheme. It is clear that we may replace $V$ by $V'$
and hence we may assume $V$ is a scheme. Since $X$ is quasi-compact,
there exist finitely many affine opens $V_i \subset V$ such that
$V' = \coprod V_i \to X$ is surjective.
After replacing $V$ by $V'$ again, we see that we may assume $V$ is affine.
Since $X$ is quasi-separated, hence reasonable, there exists an integer
$d$ bounding the degree of the geometric fibres of $V \to X$
(see Lemma \ref{lemma-bounded-fibres}).

\medskip\noindent
By induction on $d \geq 1$ we will prove the following induction
hypothesis $(H_d)$:
\begin{itemize}
\item for any quasi-compact and quasi-separated algebraic space
$X$ with finitely many irreducible components, for any $m \geq 0$,
for any quasi-compact and separated schemes $V_j$, $j = 1, \ldots, m$,
for any \'etale morphisms $\varphi_j : V_j \to X$, $j = 1, \ldots, m$
such that $d$ bounds the degree of the geometric fibres of
$\varphi_j : V_j\to X$ and
$\varphi = \coprod \varphi_j : V = \coprod V_j \to X$
is surjective, the statement of the lemma holds for $\varphi : V \to X$.
\end{itemize}
If $d = 1$, then each $\varphi_j$ is an open immersion. Hence $X$
is a scheme and the result holds with $Y = V$.
Assume $d > 1$, assume $(H_{d - 1})$ and let
$m$, $\varphi : V_j \to X$, $j = 1, \ldots, m$ be as in $(H_d)$.

\medskip\noindent
Let $\eta_1, \ldots, \eta_n \in |X|$ be the generic points of the
irreducible components of $|X|$. By
Properties of Spaces, Proposition
\ref{spaces-properties-proposition-locally-quasi-separated-open-dense-scheme}
there is an open subscheme $U \subset X$ with $\eta_1, \ldots, \eta_n \in U$.
By shrinking $U$ we may assume $U$ affine and by
Morphisms, Lemma \ref{morphisms-lemma-generically-finite}
we may assume each $\varphi_j : V_j \to X$ is finite \'etale over $U$.
Of course, we see that $U$ is quasi-compact and dense in $X$
and that $\varphi_j^{-1}(U)$ is dense in $V_j$. In particular each $V_j$
has finitely many irreducible components.

\medskip\noindent
Fix $j \in \{1, \ldots, m\}$.
As in Morphisms of Spaces, Lemma
\ref{spaces-morphisms-lemma-quasi-finite-separated-quasi-affine}
we let $Y_j$ be the normalization of $X$ in $V_j$. We obtain a factorization
$$
\xymatrix{
V_j \ar[rr] \ar[rd]_{\varphi_j} & & Y_j \ar[ld]^{\pi_j} \\
& X
}
$$
with $\pi_j$ integral and $V_j \to Y_j$ a quasi-compact open immersion. Since
$Y_j$ is the normalization of $X$ in $V_j$, we see from
Morphisms of Spaces, Lemmas
\ref{spaces-morphisms-lemma-properties-normalization} and
\ref{spaces-morphisms-lemma-normalization-in-integral}
that $\varphi_j^{-1}(U) \to \pi_j^{-1}(U)$ is an isomorphism.
Thus $\pi_j$ is finite \'etale over $U$. Observe that $V_j$
is scheme theoretically dense in $Y_j$ because $Y_j$ is the normalization
of $X$ in $V_j$ (follows from the characterization of relative normalization
in Morphisms of Spaces, Lemma
\ref{spaces-morphisms-lemma-characterize-normalization}). Since $V_j$
is quasi-compact we see that $|V_j| \subset |Y_j|$ is dense,
see Morphisms of Spaces, Section
\ref{spaces-morphisms-section-scheme-theoretic-closure}
(and especially Morphisms of Spaces, Lemma
\ref{spaces-morphisms-lemma-quasi-compact-immersion}).
It follows that $|Y_j|$ has finitely many irreducible components.
Then $V_j \times_X Y_j$ is a quasi-compact, separated scheme
(being finite over $V_j$) and
$$
V_j \times_X Y_j = V_j \amalg W_j
$$
Here the first summand is the image of $V_j \to V_j \times_X Y_j$
(which is closed by
Morphisms of Spaces, Lemma \ref{spaces-morphisms-lemma-semi-diagonal}
and open because it is \'etale as a morphism between
algebraic spaces \'etale over $Y$) and
the second summand is the (open and closed) complement.

\medskip\noindent
The \'etale morphism $W_j \to Y_j$ has geometric fibres of cardinality $< d$.
Namely, this is clear for geometric points of $V_j \subset Y_j$ by inspection.
Since $|V_j| \subset |Y_j|$ is dense, it holds for all geometric points
of $Y_j$ by Lemma \ref{lemma-quasi-compact-reasonable-stratification}
(the degree of the fibres of a quasi-compact \'etale morphism
does not go up under specialization). By $(H_{d - 1})$ applied
to $V_j \amalg W_j \to Y_j$ we find a surjective integral morphism
$Y_j' \to Y_j$ with $Y_j'$ a scheme, which Zariski locally factors
through $V_j \amalg W_j$, and which is finite \'etale over a
quasi-compact dense open $U_j \subset Y_j$. After shrinking $U$
we may and do assume that $\pi_j^{-1}(U) \subset U_j$
(we may and do choose the same $U$ for all $j$; some details omitted).

\medskip\noindent
We claim that
$$
Y = \coprod\nolimits_{j = 1, \ldots, m} Y'_j \longrightarrow X
$$
is the solution to our problem. First, this morphism is integral
as on each summand we have the composition $Y'_j \to Y \to X$
of integral morphisms (Morphisms of Spaces, Lemma
\ref{spaces-morphisms-lemma-composition-integral}). Second, this
morphism Zariski locally factors through $V = \coprod V_j$ because
we saw above that each $Y'_j \to Y_j$ factors Zariski locally through
$V_j \amalg W_j = V_j \times_X Y_j$. Finally, since both
$Y'_j \to Y_j$ and $Y_j \to X$ are finite \'etale over
$U$, so is the composition. This finishes the proof.
\end{proof}

\bibliographystyle{plain}
\bibliography{my}
\end{document}
